\chapter{Conclusão}
\label{chap:conclusoes-e-trabalhos-futuros}

Este trabalho investigou algoritmos de busca aplicados à otimização de rotas para visitação de 
ativos públicos municipais no contexto da zeladoria urbana. O objetivo central foi desenvolver e 
avaliar soluções computacionais que auxiliem gestores públicos na geração de rotas eficientes para 
equipes de manutenção que precisam visitar múltiplos pontos distribuídos geograficamente pela cidade.

\section{Atividades Realizadas}

Até o presente momento, foram desenvolvidas as seguintes etapas do trabalho:

\begin{itemize}
    \item \textbf{Revisão Sistemática da Literatura:} Conduzida seguindo o protocolo de 
    \citeonline{KITCHENHAM2009}, resultando na identificação de 8 trabalhos primários relevantes 
    sobre otimização de rotas urbanas e aplicações do \gls{TSP} em gestão pública;
    
    \item \textbf{Modelagem do Problema:} Obtenção e estruturação de grafos representando a malha 
    viária do município de Russas-CE, utilizando dados geográficos do \gls{ORS}, com grafos 
    variando de 13 a 81 nós;
    
    \item \textbf{Implementação dos Algoritmos:} Desenvolvimento em Python de cinco algoritmos de 
    busca para resolução do problema: \gls{DFS}, \gls{BFS}, \gls{UCS}, \gls{AE} com heurística 
    euclidiana e \gls{AE} com heurística haversiana;
    
    \item \textbf{Execução dos Experimentos:} Realização de testes automatizados sobre diferentes 
    instâncias do problema, variando o tamanho do grafo e a quantidade de ativos a serem visitados 
    (5 a 40 pontos);
    
    \item \textbf{Resultados Preliminares:} Obtenção de dados comparativos sobre tempo de execução 
    e distância das rotas geradas por cada algoritmo, evidenciando a superioridade dos algoritmos 
    de busca informada (\gls{UCS} e \gls{AE}) em relação aos não-informados (\gls{DFS} e \gls{BFS}).
\end{itemize}

Os resultados preliminares indicam que todos os algoritmos convergem para soluções com distâncias 
equivalentes quando há convergência, porém os algoritmos informados apresentam desempenho temporal 
substancialmente superior, com tempos de execução inferiores a 1 segundo mesmo para as instâncias 
mais complexas, enquanto os algoritmos não-informados não convergem para grafos maiores.

\section{Próximas Etapas}

Conforme estabelecido no cronograma de trabalho, as atividades previstas para conclusão do TCC são:

\begin{itemize}
    \item \textbf{Análise Aprofundada dos Resultados:} Realizar análises estatísticas mais 
    detalhadas dos dados experimentais, incluindo médias, desvios-padrão e comparações de 
    escalabilidade entre os algoritmos;
    
    \item \textbf{Desenvolvimento de Visualizações Gráficas:} Criar representações visuais das 
    rotas otimizadas sobre os mapas reais, facilitando a interpretação e validação das soluções 
    geradas;
    
    \item \textbf{Construção de Aplicação Visual:} Desenvolver uma interface gráfica demonstrativa 
    que permita visualizar o funcionamento dos algoritmos e as rotas geradas sobre a malha viária;
    
    \item \textbf{Redação Final do Documento:} Consolidar todas as seções do trabalho, incorporando 
    as análises finais, discussões aprofundadas e considerações conclusivas sobre a aplicabilidade 
    prática da solução proposta.
\end{itemize}

\section{Considerações Parciais}

Os resultados obtidos até o momento demonstram a viabilidade técnica para apoiar a
otimização de rotas na gestão pública municipal. A abordagem desenvolvida 
mostrou-se adequada para resolver problemas reais de zeladoria urbana em escalas compatíveis com 
municípios de pequeno e médio porte, oferecendo tempos de execução compatíveis com aplicações 
práticas.

Espera-se que, ao término do trabalho, seja possível fornecer subsídios técnicos concretos para 
que gestores públicos possam implementar sistemas de otimização de rotas em suas administrações, 
contribuindo para a redução de custos operacionais, economia de recursos públicos e melhoria na 
eficiência dos serviços prestados à população. A continuidade do trabalho também poderá servir 
como base para futuros desenvolvimentos na área, incluindo a avaliação de meta-heurísticas, 
incorporação de restrições adicionais e aplicação em outros contextos de roteamento urbano.


